%% kapa: exam=1 round=33 year=2022 period=1 form=A subject=real_estate subject_ko=부동산학원론 seq=8 number=88
%% source: Q-Net 감정평가사 기출문제 (cst00309) artlSeq=5209997 file=raw/1008/5209997/2022년도 제33회 감정평가사 1차시험 1교시(민법,경제학,부동산학).pdf
%% conversion: kordoc 4.24.51 -> tools/gen1.py (auto); review=proofread-sonnet; flags=none
\begin{question}{88}
산업입지이론에 관한 설명으로 옳지 않은 것은?
\begin{choices}
  \item 베버(A. Weber)는 운송비의 관점에서 특정 공장이 원료지향적인지 또는 시장지향적인지 판단하기 위해 원료지수(material index)를 사용하였다.
  \item 베버(A. Weber)의 최소비용이론에서는 노동비, 운송비, 집적이익 가운데 운송비를 최적입지 결정에 가장 우선적으로 검토한다.
  \item 뢰쉬(A. Lösch)의 최대수요이론에서는 입지분석에 있어 대상지역 내 원자재가 불균등하게 존재한다는 전제 하에, 수요가 최대가 되는 지점이 최적입지라고 본다.
  \item 아이사드(W. Isard)는 여러 입지 가운데 하나의 입지를 선정할 때 각 후보지역이 가지고 있는 비용최소 요인을 대체함으로써 최적입지가 달라질 수 있다는 대체원리(substitution principle)를 입지이론에 적용하였다.
  \item 스미스(D. Smith)의 비용수요통합이론에서는 이윤을 창출할 수 있는 공간한계 내에서는 어디든지 입지할 수 있다는 준최적입지(suboptimal location) 개념을 강조한다.
\end{choices}
\end{question}
\answer{3}
